\section{Design Checklist and Contract Skeleton}

The checklist and migration protocol spell out what a runtime must record, reserve, and recover.

\subsection{Design Checklist}

For each decision, identify the state object, the evidence needed to act, and the service property it must preserve.

\begin{table}[H]
\small
\centering
\caption{Appendix guide for translating synthesis into runtime design choices.}
\label{tab:supp-design-guide}
\begin{tabular}{@{}p{0.25\textwidth}p{0.43\textwidth}p{0.22\textwidth}@{}}
\toprule
Recurring synthesis result & Design question & Concrete artifact to look for \\
\midrule
Observability gaps dominate failure analysis & Which signal names the state object, lifecycle stage, delay semantics, and legal action? & typed telemetry and disturbance indicators \\
Local speedups disappear under composition & Which boundary is hard, which objective is soft, and who wins on conflict? & precedence-aware decision kernel \\
Lifecycle debt accumulates across updates and movement & Where are admit, place, mutate, expose/transfer, compact, and evict/reclaim made explicit? & lifecycle pipeline with bounded actuation interfaces \\
Recovery reuses the same state objects as steady state & Which invariants survive migration, replay, rebuild, or restore? & safety envelope and ownership-transfer rules \\
\bottomrule
\end{tabular}
\end{table}

\subsection{Contract Principles}

Record the state identity, observation time, and action allowed by that observation. After acting, update the record with the result and any work still pending. A later controller can then distinguish observed load from load already reserved by another action.

\subsection{Architecture Axes}

Choose a granularity that matches the legal unit of change. Place authorization where controllers can observe or reserve the relevant resources. For an object spanning workers, an embedded controller may need a shared reservation even when its own scheduling is local.

\subsection{Anti-Pattern Repairs}

For a read path that defers maintenance, measure both query performance and pending updates. For a moving shard, inject failures on either side of authority commit and verify that replay starts at the committed position.

\begin{figure*}[t]
\centering
\includegraphics[width=0.96\textwidth]{figures/pdf/antipattern_contract_map.pdf}
\caption{Anti-pattern to contract repair map.}
\label{fig:supp-antipattern-map}
\end{figure*}

\subsection{Blueprint Contract Skeleton}

The worker, link scheduler and recovery coordinator must use the same reservation for a committed transfer. Table~\ref{tab:supp-policy-skeleton} gives the order of decisions.

\begin{table*}[t]
\small
\caption{Abstract policy skeleton for unified state governance.}
\label{tab:supp-policy-skeleton}
\begin{tabular}{@{}p{0.09\textwidth}p{0.83\textwidth}@{}}
\toprule
Step & Policy skeleton \\
\midrule
1 & \texttt{observe\_signals()} and update telemetry confidence for contention, queueing, local freshness, exposure freshness, and retention pressure. \\
2 & \texttt{identify\_state\_objects()} whose boundary is currently active: hot operator shards, KV pages, retrieval indexes, replay buffers, or demotion candidates. \\
3 & Check fresh foreground demand and backpressure at both endpoints and along the path. Defer if a bound is unavailable, pressure reaches the stop threshold, or the candidate exceeds headroom after reservations. \\
4 & Compare cumulative waiting with the reference placement. Require a write pause within the service limit and enough time remaining after settlement for the move to repay its cost. \\
5 & Reserve capacity before fencing. Stop issuing batches when foreground pressure rises beyond its bound. Keep the resources and progress records required to complete or abort a committed move. \\
6 & Restart at a low rate after queues drain and backpressure remains below the lower restart threshold for a recovery interval. Revalidate reservations before the next batch. \\
\bottomrule
\end{tabular}
\end{table*}

\subsection{Boundary Case Bridges}

Appendix~C instantiates the checklist with a serving revisit, a shard handoff, and index maintenance. The three tests end when the affected request, replay, or query can be observed.

\subsection{Hardware-Conscious Evidence Anchors}

Three costs arise when state moves across hardware. The first comes from its representation: approximate execution and compressed-state systems such as Anda, TurboAttention, QServe, ATOM, VAttention, FlashInfer, SampleAttention, QuantLLM, and related low-bit cache work all show that encoding decisions reshape memory traffic, kernel regularity, and quality loss simultaneously~\cite{fang2025anda,kang2025turboattention,lin2025qserve,zhao2024atom,prabhu2025vattention,ye2025flashinfer,zhu2024sampleattention,xia2024quantllm,dong2024qhitter}. Even systems that make representation look cheaper, such as FuseMax-style lines, do so by shifting complexity into calibration, selection, or error accounting rather than eliminating it~\cite{nayak2024fusemax,kwon2020fvm,jha2025hycache}.

The second is \emph{movement debt}. Systems such as Chimera, HeterInfer, LIA, HETIS, Mercury, WaferLLM, DRAMCache, Tiertune, MoE-lightweight offloading, PhoenixOS, JIT-Serve, MLP offload, Colloid, BeyondHotness, Optimus, Tigon, and MemStrata all show that placement helps only when the runtime can jointly price compute skew, cache residency, interconnect bandwidth, and cross-device movement~\cite{huang2025chimera,chen2025heteroinfer,kim2025lia,mo2025hetis,mercury2025rms,he2025waferllm,hong2024dramcache,lee2025tiertune,cao2025moelightning,wei2025phoenixos,zhang2026jitserve,maurya2025mlpoffload,vuppalapati2024colloid,liu2025beyondhotness,feng2025optimus,huang2025tigon,zhong2024memstrata}. Storage- and I/O-adjacent systems such as AutoScratch, PFSCK, PolarDB, CXL-based designs, Toleo, and VM-control make the same point from another boundary: once state crosses devices or tiers, metadata handling, legality, and transport cost become first-order control variables rather than implementation details~\cite{fu2024autoscratch,domingo2021pfsck,cao2020polardb,ji2025cxl,dong2024toleo,zheng2024vmcu,tabatabai2024fbmm}.

The third is \emph{resilience and repair debt}. Save, ChecknRun, CXL memory-protection work, Concealing, TurboAttention-style repair-aware approximation, SampleAttention, and PFSCK show that movement and representation cannot be priced independently from failure handling, degradation tolerance, or hidden repair cost~\cite{zheng2025save,eisenman2022checknrun,zhao2025equilibria,jin2024concealing,kang2025turboattention,zhu2024sampleattention,domingo2021pfsck}. Heterogeneity-aware execution and server-placement lines such as GMLake, Flex-MoE, Kamath-style overlap, Clap, DecDec, Coruscant, Zico, MEPipe, Clone, PowerInfer, PIMBA, QFactory, and MoC reinforce the same lesson from scheduling: overlap, offload, or recomputation only help if the control overhead and recovery debt stay below the service debt they are meant to avoid~\cite{guo2024gmlake,cao2025moelightning,kamath2025podattention,park2025clap,park2025decdec,joo2025coruscant,zico21,sun2025mepipe,tian2025clone,song2024powerinfer,kim2025pimba,zhang2025qfactory,cai2025moc}.

Dynamic input pruning with cache-aware masking, near-core decompression, and CXL-centered inference expose different actuators, yet each still ties placement and movement to an explicit runtime choice over state: what may be skipped, decompressed, kept near compute, or shifted into the memory hierarchy~\cite{federici2025dip,gerogiannis2025deca,gu2025cent}. Placement, movement, and state transformation therefore remain one coupled control problem even when the local actuator changes.

Once state crosses a hardware boundary, representation, movement, and recovery costs jointly determine whether a local optimization improves service. The relative cost of each depends on the data path and the failure conditions the system must handle.
